\originalchapter{共同引力势、薄盘与暗晕模型}

本章对应 8.902 第4讲及第5讲开头的势—密度对. 对许多恒星的运动, 先不逐个计算所有配对力, 而用平滑质量密度 $\rho$ 产生共同势 $\Phi$. 一个可用模型需同时满足力的方向、Poisson 方程和边界条件. 原笔记中数处势的正负号、测试质量和 $G$ 因子互相不一致, 以下固定 $\Phi$ 为单位质量势, 以这些条件逐项复核.

\originalsection{Poisson 方程与势零点}
加速度 $\boldsymbol a=-\nabla\Phi$, 引力 Gauss 定律为 $\nabla\cdot\boldsymbol a=-4\pi G\rho$, 故 $\nabla^2\Phi=4\pi G\rho$. 使用 $\nabla^2(1/|\boldsymbol r-\boldsymbol r'|)=-4\pi\delta^{(3)}(\boldsymbol r-\boldsymbol r')$, 对有限孤立质量分布取无穷远势零,
\[
\Phi(\boldsymbol r)=-G\int\frac{\rho(\boldsymbol r')}{|\boldsymbol r-\boldsymbol r'|}\dd^3r'.
\]
负号给吸引力. $[\Phi]=L^2T^{-2}$; 测试质量的势能是 $m\Phi$, 不应把 $m$ 写入共同势本身. 自引力系统总势能为 $U=\frac12\int\rho\Phi\dd^3r$, 其中 $1/2$ 避免配对重数.

球对称时, Gauss 定律直接给
\[
\Phi'(r)=\frac{GM(<r)}{r^2},\qquad
M(<r)=4\pi\int_0^r\rho(s)s^2\dd s.
\]
于是 $\boldsymbol a=-GM(<r)\hat{\boldsymbol r}/r^2$. 反过来给定一个球对称势, 可由
\[
\rho(r)=\frac1{4\pi Gr^2}\frac{\dd}{\dd r}\left(r^2\Phi'(r)\right)
\]
恢复密度, 再检查是否非负. 给定任意看起来平滑的 $\Phi$ 并不保证它对应物理质量分布.

\originalsection{点质量与均匀球的内外连接}
点质量给 $\Phi=-GM/r$, 在 $r>0$ 满足 Laplace 方程, 质量由原点的 delta 源承担. 均匀球半径 $R$, $\rho=3M/(4\pi R^3)$, 内部质量 $M(<r)=M(r/R)^3$. 积分径向力并使 $r=R$ 的势连续, 得
\[
\Phi(r)=\begin{cases}
-\dfrac{GM}{2R^3}(3R^2-r^2),&r\le R,\\[3pt]
-\dfrac{GM}{r},&r\ge R.
\end{cases}
\]
势和一阶导数均连续, 二阶导数跨密度突变可以不同. 球内加速度正比于 $-\boldsymbol r$, 因此粒子作谐振运动, $\omega^2=GM/R^3=4\pi G\rho/3$. 球壳外质量不贡献局部力, 却改变内部势的常数; 要算逃逸速度就不能丢掉这个常数.

\begin{example}
在有限均匀球中心, $v_{\rm esc}^2=-2\Phi(0)=3GM/R$; 表面则为 $2GM/R$. 若只把内部 $\Phi\propto r^2$ 的变动部分当完整势, 会误认为中心逃逸速度为零. 逃逸问题同时依赖整个质量分布与无穷远边界条件.
\end{example}

\originalsection{Mestel 薄盘的质量在导数跳跃中}
令柱坐标为 $(R,\varphi,z)$, $s=\sqrt{R^2+z^2}$. 课程的 Mestel 势可等价写为
\[
\Phi(R,z)=v_0^2\ln\frac{s+|z|}{2r_0}.
\]
$r_0$ 只确定势的加法常数. 在 $z\ne0$ 处直接微分得到 $\nabla^2\Phi=0$. 但不能据此说总密度为零, 因为 $|z|$ 的导数在盘面跳跃. 对 $z>0$, $\partial_z\Phi=v_0^2/s$; 对 $z<0$ 则为 $-v_0^2/s$. 因而
\[
\partial_z\Phi|_{0^+}-\partial_z\Phi|_{0^-}
=\frac{2v_0^2}{R}=4\pi G\Sigma(R),\qquad
\Sigma(R)=\frac{v_0^2}{2\pi GR}.
\]
这就是 $\rho=\Sigma(R)\delta(z)$ 的分布意义 Poisson 方程. 跳跃方向给 $z>0$ 的加速度向下, 与盘的吸引相符. 在盘面 $\partial_R\Phi=v_0^2/R$, 所以 $v_c^2=R\partial_R\Phi=v_0^2$, 旋转曲线恰为平坦.

注意这与球密度 $\rho\propto r^{-2}$ 是两种不同几何的模型. 薄盘的 $M(<R)=2\pi\int_0^R\Sigma(s)s\dd s=v_0^2R/G$ 恰有同样线性尺度, 但不能据此把球壳定理套给任意盘. 两者各自以完整势满足圆轨道条件. 无限 Mestel 盘的总质量发散, 中心也奇异; 它只能在有限半径区间近似真实盘. 盘外截断会改变整体势和旋转曲线, 不能保持每条公式都精确.

\originalsection{NFW 暗晕及集中度的物理定义}
NFW 是课程用于无碰撞冷暗物质模拟的经验球形轮廓. 取尺度半径 $a$, $x=r/a$,
\[
\rho(r)=\frac{\rho_s}{x(1+x)^2}.
\]
内侧 $r\ll a$ 为 $r^{-1}$ 尖峰, 外侧为 $r^{-3}$. 质量积分用 $x/(1+x)^2=1/(1+x)-1/(1+x)^2$ 得
\[
M(<r)=4\pi\rho_sa^3\left[\ln(1+x)-\frac{x}{1+x}\right].
\]
令 $m(x)$ 为方括号, 则圆速度 $v_c^2=4\pi G\rho_sa^2m(x)/x$. 势满足无穷远零点时
\[
\Phi(r)=-4\pi G\rho_sa^2\frac{\ln(1+x)}x.
\]
微分得到 $\Phi'=GM(<r)/r^2$, 是对常数和 $G$ 因子的独立检查. 原课式漏写的 $G$ 在此恢复; 否则量纲不是速度平方. 当 $x\to0$ 势趋有限常数, 尖峰密度不必使势本身发散. 当 $x\to\infty$ 质量对数发散但势趋零, 两者也不矛盾.

定义临界密度 $\rho_c=3H^2/(8\pi G)$, $r_{200}$ 内平均密度为 $200\rho_c$, 集中度 $c_{200}=r_{200}/a$. 此处下标避免把集中度与光速 $c$ 混淆. 由两种质量表达式相等,
\[
M_{200}=\frac{4\pi}{3}200\rho_cr_{200}^3,
\qquad
\frac{\rho_s}{\rho_c}=\frac{200}{3}\frac{c_{200}^3}{m(c_{200})}.
\]
因此 $M_{200},c_{200}$ 决定模型两个参数. 模拟中的质量—集中度趋势可减少平均意义的自由度, 但同质量晕仍有形成史和散布, 不应把原笔记“同质量完全相同”当作精确唯一性定理. 以平均宇宙物质密度而非临界密度定义的 $r_{200}$ 又是另一约定, 需声明参考密度和红移.

\begin{center}
\begin{tikzpicture}
\begin{axis}[width=8cm,height=4.8cm,xmode=log,ymode=log,xlabel={$r/a$},ylabel={$\rho/\rho_s$},xmin=0.03,xmax=30,grid=major]
\addplot[structurecolor,thick,domain=0.03:30,samples=160]{1/(x*(1+x)^2)};
\end{axis}
\end{tikzpicture}
\end{center}
图直接采样轮廓, 不是对观测数据的拟合. 要把模型作为有限系统使用, 还需截断或与环境匹配; 在截断半径内力可保留同一包围质量, 势的常数及外部逃逸条件却会改变.

\originalsection{小尺度问题应比较同一种可观测量}
课程讨论核—尖峰、旋转曲线多样性、缺失卫星和过大而难以没有恒星的子晕等问题. 比较时不能把纯暗物质子晕数直接等同于已观测发光卫星数: 气体冷却、反馈、潮汐剥离及观测选择都介入. 由恒星轨道推断中央暗物质密度又受到各向异性退化, 需把动力学误差与模型预测一起比较.

这些问题激发重子反馈以及不同暗物质模型的检验, 并非其中任意一项已经单独确定了粒子性质. 本章保留2023年课程的问题结构, 不复制含权利例外的图片, 也不把课程的概述包装为2026年最新观测结论. 下一章将从给定势建立轨道, 再解释从运动学反推势时还缺哪些条件.
